\begin{corollary}[3.4.2]
\label{IV.2.3.4.2-fr}
Bajo las hipótesis generales de \textup{(3.4.1)}, supongamos que $Z$ no sea igual
a uno de los $S_i$ y que
$\operatorname{long}((\mathcal F/t\mathcal F)_z)=1$. Entonces existe uno solo
de los $S_i$ que contiene a $Z$, y para ese valor de $i$ se tiene
$\operatorname{long}(\mathcal F_{s_i})=1$~; además,
$\mathcal O_{S,z}$ es un anillo de valoración discreta cuyo $t_z$ es un
uniformizante, y $t_z$ es $\mathcal F_z$-regular.
\end{corollary}

\begin{proof}
Esto se deduce de (3.4.1), pues los dos miembros de (3.4.1.1) son entonces
iguales.
\end{proof}

\begin{proposition}[3.4.3]
\label{IV.2.3.4.3-fr}
Sean $X$ un preesquema localmente noetheriano, $t$ una sección de
$\mathcal O_X$ sobre $X$, $Y$ el subpreesquema cerrado de $X$ definido por
el ideal $t\mathcal O_X$ de $\mathcal O_X$. Sean $\mathcal F$ un
$\mathcal O_X$-módulo coherente, $T$ un ciclo primo asociado a $\mathcal F$,
$T'$ una componente irreducible de $T\cap Y$, $x$ el punto genérico de
$T'$. Supongamos que $t_x$ sea $\mathcal F_x$-regular; entonces se tiene
$x\in\operatorname{Ass}(\mathcal F/t\mathcal F)$.
\end{proposition}

\begin{proof}
Como en la demostración de (3.4.1), es posible reducirse al caso en que
$X=\operatorname{Spec}(\mathcal O_x)$; teniendo en cuenta (3.1.2), la
proposición se deduce entonces de $(\mathbf{0},16.4.6.3)$.
\end{proof}

\begin{proposition}[3.4.4]
\label{IV.2.3.4.4-fr}
Sean $X$ un preesquema localmente noetheriano, $t$ una sección de
$\mathcal O_X$ sobre $X$, $Y$ el subpreesquema cerrado de $X$ definido por
el ideal $t\mathcal O_X$ de $\mathcal O_X$. Sean $\mathcal F$ un
$\mathcal O_X$-módulo coherente, $(S_i)$ la familia de las componentes
irreducibles de $\operatorname{Supp}(\mathcal F)$. Sea $y$ un punto de $Y$
tal que $t_y$ sea $\mathcal F_y$-regular y que ninguno de los ciclos primos
asociados inmersos de $\mathcal F/t\mathcal F$ contenga a $y$. Entonces las
componentes irreducibles de $\operatorname{Supp}(\mathcal F/t\mathcal F)$ que
contienen a $y$ son exactamente las componentes irreducibles de los $S_i\cap Y$
que contienen a $y$, y los ciclos primos asociados a $\mathcal F$ que
contienen a $y$ son no inmersos.
\end{proposition}

\begin{proof}
Demostremos primero la última afirmación. Sean $T\supset T_1$ dos ciclos
primos asociados a $\mathcal F$ que contienen a $y$; si $x$ es el punto genérico
de una componente irreducible de $T\cap Y$ que contiene a $y$, $x$ es una
generización de $y$, por tanto está contenido en toda vecindad de $y$, y la
hipótesis de que $t_y$ es $\mathcal F_y$-regular implica que $t_x$ es
$\mathcal F_x$-regular $(\mathbf{0},15.2.4)$, por lo que, en virtud de (3.4.3), se tiene
$x\in\operatorname{Ass}(\mathcal F/t\mathcal F)$. Sea $x_1$ el punto
genérico de una componente irreducible de $T_1\cap Y$ que contiene a $y$, y sea
$x$ el punto genérico de una componente irreducible de $T\cap Y$ que contiene a
$x_1$; se sigue de lo anterior que $x_1$ y $x$ pertenecen ambos a
$\operatorname{Ass}(\mathcal F/t\mathcal F)$, y como
$x_1\in\overline{\{x\}}$, la hipótesis implica que $x_1=x$. Designemos también
por $T$ y $T_1$ los subpreesquemas cerrados íntegros de $X$ que tienen
respectivamente a $T$ y $T_1$ como espacios subyacentes, y pongamos
$A=\mathcal O_{T,x}$, $A_1=\mathcal O_{T_1,x}$; se tiene, pues,
$A_1=A/\mathfrak p$, donde $\mathfrak p$ es un ideal primo de $A$. Por
definición de $x$ y $x_1$, $A/tA$ y $A_1/tA_1$ son anillos \emph{artinianos};
por otra parte, se ha visto más arriba que $t_x$ es $\mathcal F_x$-regular, de
modo que $x$ no puede pertenecer a $\operatorname{Ass}(\mathcal F)$, y por
consiguiente $A_1$ no es artiniano. Por tanto, $\dim A=\dim A_1=1$; pero esto
implica $\mathfrak p=0$ y $A=A_1$ $(\mathbf{0},16.1.2.2)$; como
$\operatorname{Spec}(A)$ y $\operatorname{Spec}(A_1)$ son respectivamente
densos en $T$ y $T_1$, se tiene efectivamente $T=T_1$.

Los $S_i$ que contienen a $y$ son, por tanto, todos los ciclos primos asociados
a $\mathcal F$ que contienen a $y$; si $x$ es el punto genérico de una
componente irreducible de $S_i\cap Y$ que contiene a $y$, se deduce de nuevo
de $(\mathbf{0},15.2.4)$ que $t_x$ es $\mathcal F_x$-regular, por lo que, en virtud de
(3.4.3), $x\in\operatorname{Ass}(\mathcal F/t\mathcal F)$; esto demuestra la
primera afirmación de (3.4.4).
\end{proof}

\begin{proposition}[3.4.5]
\label{IV.2.3.4.5-fr}
Sean $X$ un preesquema localmente noetheriano, $t$ una sección de
$\mathcal O_X$ sobre $X$, $Y$ el subpreesquema cerrado de $X$ definido por
el ideal $t\mathcal O_X$ de $\mathcal O_X$. Sean $\mathcal F$

\oldpage[IV-2]{50}

un $\mathcal O_X$-módulo coherente, $y$ un punto de $Y$; supongamos que $t_y$
sea $\mathcal F_y$-regular y que $\mathcal F/t\mathcal F$ sea íntegro en el
punto $y$ \textup{(3.2.4)}. Entonces $\mathcal F$ es íntegro en el punto $y$.
\end{proposition}

\begin{proof}
Teniendo en cuenta (3.4.4), basta demostrar que $y$ está contenido en \emph{una
sola} componente irreducible de $\operatorname{Supp}(\mathcal F)$, y
que si $s$ es el punto genérico de esta componente, se tiene
$\operatorname{long}(\mathcal F_s)=1$. Ahora bien, por hipótesis, $y$ no
pertenece más que a una sola componente irreducible de
$\operatorname{Supp}(\mathcal F/t\mathcal F)$, y si $z$ es el punto
genérico de esta componente, se tiene
$\operatorname{long}((\mathcal F/t\mathcal F)_z)=1$~; la conclusión se deduce
pues de (3.4.2).
\end{proof}
